\frfilename{mt1a1.tex}
\versiondate{5.11.03}
\copyrightdate{1996}

\def\chaptername{Lampiran}
\def\sectionname{Teori himpunan}

\newsection{1A1}

Pada 111E-111F saya membahas secara ringkas himpunan `terhitung'. Pendekatan
di sana mengikuti jalur yang tampaknya paling singkat menuju fakta-fakta yang
segera diperlukan, dan mungkin sudah selayaknya saya menunjukkan di sini jalur
yang lebih lazim. Saya memanfaatkan kesempatan ini untuk mencantumkan beberapa
notasi yang menurut saya praktis, tetapi tidak digunakan secara universal.

\leader{1A1A}{Notasi kurung siku}\cmmnt{ Saya menggunakan kurung siku
$[$ dan $]$ dalam berbagai cara; saya berharap konteks selalu memperjelas
penafsiran yang dimaksud.

\header{1A1Aa}}{\bf (a)} Untuk $a$, $b\in\Bbb R$, saya menulis

\Centerline{$[a,b]=\{x:a\le x\le b\}$,
\quad$\ooint{a,b}=\{x:a<x<b\}$,}

\Centerline{$\coint{a,b}=\{x:a\le x<b\}$,
\quad$\ocint{a,b}=\{x:a<x\le b\}$.}

\noindent Ketika rumus-rumus ini muncul, wajar jika kita langsung
menyimpulkan bahwa $a<b$; tetapi sesekali berguna untuk menafsirkannya ketika
$b\le a$, dan dalam hal itu\cmmnt{ saya mengikuti rumus-rumus di atas secara
harfiah, sehingga}

\Centerline{$[a,a]=\{a\}$,
\quad$\ooint{a,a}=\coint{a,a}=\ocint{a,a}=\emptyset$,}

\Centerline{$[a,b]=\ooint{a,b}=\coint{a,b}=\ocint{a,b}=\emptyset$ jika
$b<a$.}

\spheader 1A1Ab\cmmnt{ Kita dapat menafsirkan rumus-rumus tersebut dengan
$a$ atau $b$ yang tak hingga;
sebagai contoh,}

\Centerline{$\ooint{-\infty,b}=\{x:x<b\}$,
\quad$\ooint{a,\infty}=\{x:a<x\}$,
\quad$\ooint{-\infty,\infty}=\Bbb R$,}

\Centerline{$\coint{a,\infty}=\{x:x\ge a\}$,
\quad$\ocint{-\infty,b}=\{x:x\le b\}$,}

\cmmnt{\noindent dan bahkan}

\Centerline{$[0,\infty]=\{x:x\in\Bbb R,\,x\ge 0\}\cup\{\infty\}$,
\quad$[-\infty,\infty]=\Bbb R\cup\{-\infty,\infty\}$.}

\cmmnt{\spheader 1A1Ac Dengan kehati-hatian tertentu --- karena masih ada
pilihan-pilihan yang harus dibuat, yang lebih aman dinyatakan secara eksplisit
ketika diperlukan --- kita dapat menggunakan rumus serupa untuk `interval'
dalam ruang multidimensi $\BbbR^r$; lihat, misalnya, 115A atau 136D; dan bahkan
dalam himpunan terurut parsial umum, walaupun himpunan-himpunan ini belum akan
penting bagi kita sebelum Jilid 3.
}

\cmmnt{\spheader 1A1Ad Barangkali saya perlu menjelaskan pilihan saya untuk
menggunakan $\ooint{a,b}$, $\coint{a,b}$ alih-alih $(a,b)$, $[a,b)$, yang
keduanya lebih lazim dan lebih sedap dipandang. Alasan pertama hanyalah karena
rumus

\Centerline{$(1,2)\in\ooint{0,2}\times\ooint{1,3}$}

\noindent lebih mudah dipahami daripada bentuk padanannya. Secara umum,
pilihan ini menghasilkan keseimbangan yang sedikit lebih baik antara banyaknya
kemunculan $($ dan $[$, bahkan dengan memperhitungkan penggunaan lebih lanjut
dari $[\ldots]$ yang akan segera saya jelaskan.
}

\leader{1A1B}{Citra langsung dan citra balik}\cmmnt{ Sekarang saya
menjelaskan penggunaan kurung siku yang sama sekali berbeda, yang termasuk
dalam teori himpunan abstrak, bukan teori sistem bilangan real.

\header{1A1Ba}}{\bf (a)} Jika $f$ adalah suatu fungsi dan $A$ suatu himpunan,
saya menulis

\Centerline{$f[A]=\{f(x):x\in A\cap\dom f\}$}

\noindent untuk {\bf citra langsung} $A$ oleh $f$. \cmmnt{Perhatikan bahwa
walaupun $A$ sering kali merupakan subhimpunan dari domain $f$, hal ini tidak
diasumsikan.}

\spheader 1A1Bb Jika $f$ adalah suatu fungsi dan $B$ suatu himpunan, saya
menulis

\Centerline{$f^{-1}[B]=\{x:x\in\dom f,\,f(x)\in B\}$}

\noindent untuk {\bf citra balik} $B$ oleh $f$. \cmmnt{Kali ini penting untuk
memperhatikan bahwa tidak ada anggapan bahwa $f$ injektif, atau bahwa $f^{-1}$
merupakan suatu fungsi; rumus $f^{-1}[\,\,]$ diberi makna yang tidak bergantung
pada makna apa pun dari ungkapan $f^{-1}$. Namun, mudah dilihat bahwa ketika
$f$ injektif, sehingga kita mempunyai fungsi invers sejati $f^{-1}$ (yang
didefinisikan pada himpunan nilai $f$, yaitu $f[\dom f]$), maka $f^{-1}[B]$,
sebagaimana didefinisikan di sini, sama dengan penafsirannya menurut (a).}

\spheader 1A1Bc Sekarang andaikan bahwa $R$ merupakan suatu relasi, yaitu
suatu himpunan pasangan terurut, dan $A$, $B$ adalah himpunan. Maka\cmmnt{ saya
menulis}

\Centerline{$R[A]=\{y:\,\exists\,x\in A$ sedemikian sehingga $(x,y)\in R\}$,}

\Centerline{$R^{-1}[B]=\{x:\,\exists\,y\in B$ sedemikian sehingga $(x,y)\in R\}$.}

\cmmnt{\noindent Jika kita menulis

\Centerline{$R^{-1}=\{(y,x):(x,y)\in R\}$,}

\noindent maka kita memperoleh penafsiran lain untuk $R^{-1}[B]$ yang sama
dengan penafsiran yang baru diberikan. Selain itu, jika $R$ merupakan grafik
suatu fungsi $f$, yaitu jika untuk setiap $x$ terdapat paling banyak satu $y$
sedemikian sehingga $(x,y)\in R$, maka rumus-rumus di sini sama dengan
rumus-rumus dalam (a)-(b) di atas.
}%end of comment

\cmmnt{\spheader 1A1Bd (Bagian berikut ditujukan khusus kepada para pembaca
yang telah diajari untuk membedakan kata `himpunan' dan `kelas'.) Saya telah
menggunakan kata `himpunan' lebih dari sekali di atas. Namun, hal itu semata-
mata demi keluwesan ungkapan. Rumus-rumus yang sama dapat digunakan dengan
kelas-kelas sembarang, walaupun dalam beberapa teori himpunan ungkapan yang
terlibat mungkin tidak diakui sebagai `suku' dalam arti teknis.
}%end of comment

\cmmnt{
\leader{1A1C}{Himpunan terhitung} Dalam 111Fa saya mendefinisikan `himpunan
terhitung' sebagai berikut: suatu himpunan $K$ terhitung jika himpunan itu
kosong atau terdapat suatu fungsi surjektif dari $\Bbb N$ ke $K$. Rumusan yang
lebih lazim mengatakan bahwa suatu himpunan $K$ terhitung jika dan hanya jika
himpunan itu berhingga atau terdapat suatu bijeksi antara $\Bbb N$ dan $K$.
Jadi, saya perlu segera memeriksa bahwa kedua rumusan ini setara.
}%end of comment

\leader{1A1D}{Proposisi} Misalkan $K$ suatu himpunan. Maka pernyataan-
pernyataan berikut ekuivalen:

(i) $K$ kosong atau terdapat suatu surjeksi dari $\Bbb N$ ke $K$;

(ii) $K$ berhingga atau terdapat suatu bijeksi antara $\Bbb N$ dan $K$;

(iii) terdapat suatu injeksi dari $K$ ke $\Bbb N$.

\proof{{\bf (a)(i)$\Rightarrow$(iii)} Andaikan (i). Jika $K$ kosong, maka
fungsi kosong merupakan suatu injeksi dari $K$ ke $\Bbb N$. Jika tidak,
terdapat suatu surjeksi $\phi:\Bbb N\to K$.
Sekarang, untuk setiap $k\in K$, tetapkan

\Centerline{$\psi(k)=\min\{n:n\in\Bbb N$, $\phi(n)=k\}$;}

\noindent nilai ini selalu terdefinisi dengan baik karena $\phi$ surjektif,
sehingga $\{n:\phi(n)=k\}$ tidak pernah kosong dan mesti mempunyai anggota
terkecil. Karena $\phi\psi(k)=k$ untuk setiap $k$, $\psi$ mesti injektif, dan
karena itu merupakan injeksi yang diperlukan dari $K$ ke $\Bbb N$.

\medskip

{\bf (b)(iii)$\Rightarrow$(ii)} Andaikan (iii); misalkan
$\psi:K\to\Bbb N$ suatu injeksi, dan tetapkan
$A=\psi[K]\subseteq\Bbb N$. Maka $\psi$ merupakan suatu bijeksi antara $K$
dan $A$. Jika $K$ berhingga, tentu saja (ii) terpenuhi. Jika tidak, $A$ juga
mesti tak hingga. Definisikan $\phi:A\to\Bbb N$ dengan menetapkan

\Centerline{$\phi(m)=\#(\{i:i\in A,\,i<m\})$,}

\noindent yaitu banyaknya anggota $A$ yang kurang dari $m$, untuk setiap
$m\in A$; dengan demikian, $\phi(m)$ adalah posisi $m$ jika anggota-anggota
$A$ diurutkan dari yang terkecil, dimulai dengan $0$ untuk anggota terkecil
$A$. Maka $\phi:A\to\Bbb N$ merupakan suatu bijeksi, karena $A$ tak hingga,
dan $\phi\psi:K\to\Bbb N$ merupakan suatu bijeksi.

\medskip

{\bf (c)(ii)$\Rightarrow$(i)} Jika $K$ kosong, jelas himpunan itu memenuhi (i). Jika
$K$ berhingga dan tak kosong, daftarkan anggota-anggotanya sebagai
$k_0,\ldots,k_n$; sekarang tetapkan $\phi(i)=k_i$ untuk $i\le n$, dan $k_0$
untuk $i>n$, sehingga diperoleh suatu surjeksi $\phi:\Bbb N\to K$. Jika $K$
tak hingga, terdapat suatu bijeksi dari $\Bbb N$ ke $K$, yang tentu saja juga
merupakan suatu surjeksi dari $\Bbb N$ ke $K$. Jadi, (i) benar dalam semua
kasus.
}%end of proof of 1A1D

\cmmnt{\medskip

\noindent{\bf Catatan} Dalam bukti di atas saya menyebut `fungsi kosong'.
Fungsi ini mempunyai domain $\emptyset$; setelah hal itu dinyatakan, aturan
apa pun, atau tanpa aturan sama sekali, untuk menghitung fungsi tersebut akan
memberikan hasil yang sama karena aturan itu tidak akan pernah diterapkan.
Dengan menelaah perasaan Anda terhadap konstruksi ini, Anda dapat mengetahui
sesuatu tentang sikap dasar Anda terhadap matematika. Anda mungkin merasa
bahwa konstruksi ini merupakan sesuatu yang dibuat-buat dan tidak relevan,
atau Anda mungkin merasa bahwa konstruksi ini sama perlunya dengan bilangan
$0$. Kedua perasaan itu sepenuhnya sah, dan matematikawan yang matang akan
berganti-ganti di antara keduanya; tetapi harus saya katakan bahwa saya
sendiri lebih sering condong pada perasaan yang kedua daripada yang pertama,
dan bahwa ketika saya mengatakan `fungsi' dalam risalah ini, biasanya
kemungkinan fungsi kosong tetap ada dalam benak saya.
}%end of comment

\leader{1A1E}{Sifat-sifat himpunan terhitung}\cmmnt{ Izinkan saya
merangkum kembali sifat-sifat dasar himpunan terhitung:

\medskip

}{\bf (a)} Jika $K$ terhitung dan $\phi:K\to L$ merupakan suatu
surjeksi, maka $L$ terhitung. \prooflet{\Prf\ Jika $K$ kosong maka $L$ juga
kosong. Jika tidak, terdapat suatu surjeksi $\psi:\Bbb N\to K$, sehingga
$\phi\psi$ merupakan suatu surjeksi dari $\Bbb N$ ke $L$, dan $L$
terhitung. \Qed}

\header{1A1Eb}{\bf (b)} Jika $K$ terhitung dan $\phi:L\to K$ merupakan
suatu injeksi, maka $L$ terhitung. \prooflet{\Prf\ Menurut 1A1D(iii), terdapat
suatu injeksi $\psi:K\to\Bbb N$; sekarang $\psi\phi:L\to\Bbb N$ injektif,
sehingga $L$ terhitung. \Qed}

\header{1A1Ec}{\bf (c)} Khususnya, setiap subhimpunan dari suatu himpunan
terhitung adalah terhitung\cmmnt{ (seperti dalam 111F(b-i))}.

\header{1A1Ed}{\bf (d)} Hasil kali Kartesius dari sejumlah berhingga himpunan
terhitung adalah terhitung\cmmnt{ (111Fb(iii)-(iv))}.

\header{1A1Ee}{\bf (e)} $\Bbb Z$ terhitung. \prooflet{\Prf\ Pemetaan
$(m,n)\mapsto m-n:\Bbb N\times\Bbb N\to\Bbb Z$ bersifat surjektif. \Qed}

\header{1A1Ef}{\bf (f)} $\Bbb Q$ terhitung. \prooflet{\Prf\ Pemetaan
$(m,n)\mapsto \bover{m}{n+1}:\Bbb Z\times\Bbb N\to\Bbb Q$ bersifat surjektif.
\Qed}

\leader{1A1F}{}\cmmnt{ Ada satu lagi sifat mendasar yang patut dibedakan
dari sifat-sifat ini karena sifat tersebut bergantung pada argumen yang
sedikit lebih mendalam.

\medskip

\noindent}{\bf Teorema} Jika $\Cal K$ merupakan suatu keluarga terhitung yang
terdiri atas himpunan-himpunan terhitung, maka

\Centerline{$\bigcup\Cal K=\{x:\,\exists\,K\in\Cal K,\, x\in K\}$}

\noindent terhitung.

\proof{ Tetapkan

\Centerline{$\Cal K'=\Cal K\setminus\{\emptyset\}
=\{K:K\in\Cal K,\,K\ne\emptyset\}$;}

\noindent maka $\Cal K'\subseteq\Cal K$, sehingga himpunan ini terhitung, dan
$\bigcup\Cal K'=\bigcup\Cal K$. Jika $\Cal K'=\emptyset$, maka

\Centerline{$\bigcup\Cal K=\bigcup\Cal K'=\emptyset$}

\noindent jelas terhitung. Jika tidak, misalkan
$m\mapsto K_m:\Bbb N\to\Cal K'$ suatu surjeksi. Untuk setiap
$m\in\Bbb N$, $K_m$ merupakan suatu himpunan terhitung tak kosong, sehingga
terdapat suatu surjeksi $n\mapsto k_{mn}:\Bbb N\to K_m$. Sekarang
$(m,n)\mapsto k_{mn}:\Bbb N\times\Bbb N\to\bigcup\Cal K$ merupakan suatu
surjeksi (jika $k\in\bigcup\Cal K$, terdapat $K\in\Cal K'$ sedemikian sehingga
$k\in K$; terdapat $m\in\Bbb N$ sedemikian sehingga $K=K_m$; terdapat
$n\in\Bbb N$ sedemikian sehingga $k=k_{mn}$). Jadi, $\bigcup\Cal K$
terhitung, sebagaimana diperlukan.
}%end of proof of 1A1F

\cmmnt{
\leader{*1A1G}{Catatan} Saya memisahkan hasil ini dari fakta-fakta
`elementer' dalam 1A1E, antara lain karena hasil ini menggunakan asas argumen
yang berbeda dari semua asas yang diperlukan untuk pembahasan sebelumnya. Di
tengah bukti saya menulis `sehingga terdapat suatu surjeksi
$n\mapsto k_{mn}:\Bbb N\to K_m$'. Keberadaan suatu surjeksi dari $\Bbb N$ ke
$K_m$ memang mengikuti pernyataan tepat sebelumnya, yaitu `$K_m$ merupakan
suatu himpunan terhitung tak kosong'. Langkah terselubungnya terletak pada
pemberian nama secara langsung kepada surjeksi semacam itu sebagai
`$n\mapsto k_{mn}$'. Tentu mungkin ada banyak surjeksi dari $\Bbb N$ ke $K_m$
--- bahkan, pada umumnya, akan ada tak terhitung banyaknya --- dan pada
hakikatnya yang saya lakukan di sini ialah memilih salah satunya secara
sembarang. Pilihan itu harus sembarang karena saya bekerja dalam konteks yang
sepenuhnya abstrak, dan meskipun dalam setiap penerapan khusus teorema ini
mungkin terdapat suatu surjeksi alami untuk digunakan, saya tidak mempunyai
cara untuk meramalkan pendekatan apa, kalau pun ada, yang dapat menyediakan
kriteria untuk menentukan satu fungsi tertentu di sini. Sejak setidaknya
zaman Euklides, salah satu metode dasar argumen matematika ialah kesediaan kita
untuk memberi nama kepada suatu objek, sebuah `titik umum' atau `bilangan
sembarang', tanpa menentukan secara persis objek mana yang kita namai. Namun,
di sini saya secara serentak memilih tak hingga banyak objek, masing-masing
secara sembarang dari suatu himpunan tertentu. Ini merupakan
penggunaan {\bf Aksioma Pilihan}.

Saya tidak ingat pernah mendapati mahasiswa yang mengkritik argumen berbentuk
seperti argumen dalam 1A1F dengan alasan bahwa argumen itu menggunakan suatu
asas baru yang mungkin tidak sah; saya yakin tidak pernah terpikir oleh saya
bahwa ada sesuatu yang patut dipersoalkan dalam kasus-kasus semacam ini sampai
seseorang menunjukkannya. Jika Anda merasa bahwa pembahasan semacam ini tidak
relevan dengan minat matematika Anda sendiri, Anda tentu dapat melewatinya,
setidaknya sampai Anda mencapai Jilid 5. Sistem-sistem matematika yang di
dalamnya aksioma pilihan tidak berlaku telah dipelajari; sistem-sistem itu
sangat menarik, tetapi sejauh ini tetap berada di pinggiran bidang ini.
Sistem-sistem yang di dalamnya aksioma pilihan gagal sedemikian rupa sehingga
gabungan terhitung dari himpunan-himpunan terhitung kadang-kadang tidak
terhitung mempunyai watak tersendiri, dan khususnya teori ukuran Lebesgue
berubah secara mendasar; saya akan membahas kemungkinan ini dalam Bab 56
Jilid 5.

Untuk komentar singkat mengenai cara-cara lain menggunakan aksioma pilihan,
lihat 134C.
}%end of comment

\leader{1A1H}{Beberapa himpunan tak terhitung}\cmmnt{ Tentu saja tidak
semua himpunan terhitung. Dalam 114G/115G saya menyatakan bahwa semua
subhimpunan terhitung dari ruang Euklides terabaikan terhadap ukuran Lebesgue;
akibatnya, setiap himpunan yang tidak terabaikan --- misalnya setiap interval
tak-sepele --- mesti tak terhitung. Namun, barangkali akan membantu jika saya
menyajikan di sini argumen elementer yang menunjukkan bahwa $\Bbb R$ dan
$\Cal P\Bbb N$ tidak terhitung.

\medskip

}{\bf (a)} Tidak ada surjeksi dari $\Bbb N$ ke $\Bbb R$.
\prooflet{\Prf\ Misalkan $n\mapsto a_n:\Bbb N\to\Bbb R$ suatu fungsi
sembarang. Untuk setiap $n\in\Bbb N$, nyatakan $a_n$ dalam bentuk desimal
sebagai

\Centerline{$a_n=k_n + 0\cdot\epsilon_{n1}\epsilon_{n2}\ldots
=k_n+\sum_{i=1}^{\infty}10^{-i}\epsilon_{ni}$,}

\noindent dengan $k_n\in\Bbb Z$ bilangan bulat terbesar yang tidak melebihi
$a_n$, dan setiap $\epsilon_{ni}$ suatu bilangan bulat antara $0$ dan $9$;
sebagai konvensi, jika $a_n$ kebetulan mempunyai bentuk desimal berhingga,
gunakan ekspansi berhingga tersebut, sehingga $\epsilon_{ni}$ akhirnya selalu $0$, bukan
akhirnya selalu $9$.

Sekarang definisikan $\epsilon_i$, untuk $i\ge 1$, dengan menetapkan bahwa

$$\eqalign{\epsilon_i&=6\text{ jika }\epsilon_{ii}<6,\cr
&=5\text{ jika }\epsilon_{ii}\ge 6.\cr}$$

\noindent Tinjau $a=k_0+1+\sum_{i=1}^{\infty}10^{-i}\epsilon_i$, sehingga
$a=k_0+1+0\cdot\epsilon_1\epsilon_2\ldots$ dalam bentuk desimal. Saya
menyatakan bahwa $a\ne a_n$ untuk setiap $n$. Tentu saja $a\ne a_0$ karena
$a_0<k_0+1\le a$. Jika $n\ge 1$, maka $\epsilon_n\ne\epsilon_{nn}$; karena
tidak ada $\epsilon_i$ yang sama dengan $0$ ataupun $9$, tidak terdapat
ekspansi desimal lain untuk $a$, sehingga ekspansi
$a_n=k_n+0\cdot\epsilon_{n1}\epsilon_{n2}\ldots$ tidak dapat merepresentasikan
$a$, dan $a\ne a_n$.

Dengan demikian saya telah membangun suatu bilangan real yang tidak terdapat
dalam daftar $a_0,a_1,\ldots$. Karena $\sequencen{a_n}$ sembarang, tidak ada
surjeksi dari $\Bbb N$ ke $\Bbb R$. \Qed}

Jadi, $\Bbb R$ tak terhitung.

\header{1A1Hb}{\bf (b)} Tidak ada surjeksi dari $\Bbb N$ ke himpunan kuasanya
$\Cal P\Bbb N$. \prooflet{\Prf\ Misalkan
$n\mapsto A_n:\Bbb N\to\Cal P\Bbb N$ suatu fungsi sembarang. Tetapkan

\Centerline{$A=\{n:n\in\Bbb N,\,n\notin A_n\}$.}

\noindent Jika $n\in\Bbb N$, maka

\inset{{\it entah} $n\in A_n$, yang mengakibatkan $n\notin A$,}

\inset{{\it atau} $n\notin A_n$, yang mengakibatkan $n\in A$.}

\noindent Jadi, dalam kedua kasus kita mempunyai $n\in A\symmdiff A_n$,
sehingga $A\ne A_n$. Karena $n$ sembarang,
$A\notin\{A_n:n\in\Bbb N\}$ dan $n\mapsto A_n$ bukan suatu surjeksi. Karena
$\sequencen{A_n}$ sembarang, tidak ada surjeksi dari $\Bbb N$ ke
$\Cal P\Bbb N$. \Qed}

Jadi, $\Cal P\Bbb N$\cmmnt{ juga} tak terhitung.

\cmmnt{
\leader{1A1I}{Catatan} Sesungguhnya terdapat suatu bijeksi antara $\Bbb R$
dan $\Cal P\Bbb N$ (2A1Ha); dengan demikian, ketakterhitungan keduanya dapat
dibuktikan hanya dengan salah satu dari argumen di atas.
}

\leader{1A1J}{Notasi}\cmmnt{ Untuk memperjelas, saya nyatakan di sini bahwa}
Saya akan mengatakan bahwa suatu keluarga himpunan $\Cal A$ merupakan suatu
{\bf partisi} dari suatu himpunan $X$ jika $\Cal A$ merupakan penutup
saling lepas dari $X$, yaitu $X=\bigcup\Cal A$ dan
$A\cap A'=\emptyset$ untuk semua $A$, $A'\in\Cal A$ yang berlainan\cmmnt{;
khususnya, himpunan kosong boleh termasuk atau tidak termasuk dalam $\Cal A$}.
Demikian pula, suatu keluarga berindeks $\familyiI{A_i}$ merupakan suatu
{\bf partisi} dari $X$ jika $\bigcup_{i\in I}A_i=X$ dan
$A_i\cap A_j=\emptyset$ untuk semua $i$, $j\in I$ yang berlainan\cmmnt{;
sekali lagi, satu atau beberapa $A_i$ boleh kosong}.

\exercises{}

\cmmnt{\Notesheader{1A1} Gagasan-gagasan dalam
1A1C-1A1I %1A1C 1A1D 1A1E 1A1F 1A1H 1A1I
pada dasarnya berasal dari G.F.Cantor. Konsep-konsep ini mendasar bagi teori
himpunan modern, dan bahkan bagi bagian-bagian yang sangat luas dari
matematika murni modern. Catatan-catatan di atas baru sedikit sekali
menggambarkan kesuburan luar biasa dari gagasan-gagasan ini, yang memerlukan
buku tersendiri agar dapat diuraikan secara semestinya; satu-satunya tujuan
saya di sini ialah mencoba menjelaskan bagian-bagian kecil dari bidang ini
yang diperlukan dalam jilid sekarang. Dalam jilid-jilid selanjutnya saya akan
menyajikan hasil-hasil yang menggunakan gagasan-gagasan yang jauh lebih maju, yang
akan saya bahas dalam lampiran-lampiran jilid tersebut.
}%end of notes

\discrpage
